\section{Proposed Work and Conclusion}
\label{sec:proposed_work}

Based on the analysis of these publications, we propose implementing a \textbf{Mask R-CNN with a ResNet-101-FPN backbone}. We will utilize the \textbf{Chula-Parasite-Egg-11} dataset \cite{anantrasirichai2022icip} for training.

The analysis confirms that while detection (bounding boxes) is well-studied \cite{viet2019parasite}, instance segmentation (masks) is the necessary next step for clinical precision \cite{cringoli2024automating}. By incorporating attention mechanisms \cite{wang2024automated} and robust augmentation strategies identified in the literature, our project aims to build a tool capable of reliable automated diagnosis in resource-constrained settings.

\subsection{Addressing the Segmentation Gap: From Boxes to Masks}
A critical challenge in this study is the discrepancy between the model's requirements and the available data. Mask R-CNN is fully supervised and requires pixel-level instance masks for training, yet the Chula-Parasite-Egg-11 dataset provides only bounding box annotations. While our current implementation utilizes placeholder masks (filling the bounding box) to enable training, this approach prevents the model from learning the true morphological features of the parasites.

To overcome this limitation in future iterations, we propose a two-stage \textbf{Pseudo-Mask Generation Pipeline}:

\subsubsection{Stage 1: Classical Segmentation Initialization}
We will generate initial "weak" masks using classical computer vision techniques within the provided bounding boxes. Methods such as \textbf{Otsu's Thresholding} or \textbf{GrabCut} algorithm can effectively separate the parasite egg from the immediate background based on intensity and texture gradients. This provides a significantly better training signal than a rectangular box.

\subsubsection{Stage 2: GAN-Based Refinement}
To further refine these weak masks, we propose utilizing a \textbf{Generative Adversarial Network (GAN)}. A conditional GAN (cGAN), such as Pix2Pix, can be trained to translate the rough, classical segmentation outputs into high-fidelity, realistic binary masks. The GAN would effectively "learn" the shape priors of parasite eggs (e.g., the smooth oval shape of \textit{Ascaris} vs. the barrel shape of \textit{Trichuris}), correcting the irregularities and noise produced by the classical methods. This hybrid approach leverages the robustness of classical methods with the generative power of deep learning to synthesize high-quality segmentation data from bounding box annotations.
